\section*{Preface}
\addcontentsline{toc}{section}{Preface}

The idea of this book arose in a conversation with H.A. Bethe, who remarked that a little book about modern field theory which contained only Memorable Results would be a Good Thing. In the field of historical research this approach led to the publication of a treatise
\begingroup
\renewcommand{\thefootnote}{\ensuremath{\dagger}}%
\footnote{W.C. Sellar and R.J. Yeatman, \emph{1066 and All That}, Dutton, New York, 1931.}
\endgroup
which has become a standard text for serious students. Although it is often dangerous to use the tried and true methods of one subject in another field of research, the application to physics of the principles of that book has led to at least one good result: we have eliminated all theorems whose proofs are non-existent.

R.F. Streater\par
A.S. Wightman\par
November 1963

In the 1978 edition of this book, we added Appendix A, which outlined three significant developments in the general theory of quantized fields that had occurred since the appearance of the first edition: constructive quantum field theory, the theory of local algebras, and the theory of superselection rules.

Neither of these first two editions contained an account of the Haag-Ruelle collision theory. For that the reader of the first edition was referred to the then forthcoming excellent book by R. Jost (\emph{The General Theory of Quantized Fields}, American Mathematical Society, 1965). For the 1978 edition this reference was supplemented by a reference to the treatise \emph{Introduction to Axiomatic Quantum Field Theory} by N. Bogolubov, A. Logunov and I. Todorov, W.A. Benjamin (1975). For the present Princeton University Press reissue (2000), we want to add a
third recommendation to this list: H. Araki's \emph{Mathematical Theory of Quantum Fields}, Oxford University Press, 2000. In other respects the Princeton University Press edition differs from the 1978 edition mainly in the correction of a few misprints.

R.F. Streater\par
A.S. Wightman
